% Draft of the supplement subsection for study IG.  To use it once experiments/results/ig_summary.json exists:
%   cp docs/paper/ig_section.draft.tex docs/paper/ig_section.tex  (then edit the "Reading" paragraph to what the numbers show).  The include is already in supplement.tex (\IfFileExists).
\subsection{Study IG: The Request Parser on Fresh Operator Texts}\label{sec:moved-ig}
Every bound of the paper is relative to the requested supply, which the parser reads from the operator's words; its prompt was tuned on 40 development utterances and its model chosen on a 60-text held-out set, so no clean test remained. Study IG (registered and stamped before any inference) runs the frozen parser (Qwen3.5-9B, prompt i35) on \IGTexts{} fresh texts written by hand by the assistant that built it: polite, casual, terse, speech-recognition style, English, mixed, indirect needs and negations (accept expected), and two supplies, vague, wrong shelter, unsupported supply, irrelevant and wording traps (refusal expected).
\emph{Claims.} IG1: at most 2 unsafe outcomes (a wrong supply, or an accepted text that should be refused): \IGUnsafe{} (Wilson upper bound \IGUnsafeUpper): \IGOneHolds. IG2: at least 80\% of the legitimate texts accepted with the right supply: \IGAccepted{} (\IGAcceptedRate): \IGTwoHolds. Refusals correct: \IGRefused. Predictions: speech-recognition style \IGAsr{} \IGPredOneHolds; English \IGEnglish{} \IGPredTwoHolds. Other categories (accepted correctly): polite \IGPolite, casual \IGCasual, terse \IGTerse, indirect \IGIndirect, negation \IGNegation.
Japanese multimodal benchmarks compare models with English on translated subjects~\cite{onohara2024jmmmu} and measure the gap of Japanese VLMs to proprietary ones~\cite{inoue2024heronbench}, and multilingual LLMs answer translated queries inconsistently, especially for some writing systems~\cite{gupta2025foundintranslation}, so IG reports Japanese and English texts separately. Spoken commands are already parsed into ordered pickup and delivery steps by LLMs on robots~\cite{srivastava2024speechnav}, so the parser's errors are a deployment question, not only a benchmark one.
\emph{Reading.} IG1 fails, and the five unsafe texts show one cause and one oddity (all five are listed verbatim in \texttt{ig\_summary.json}). Four name a destination the mission does not know (``to the community hall'', ``to the city hall'', ``to shelter 1'', ``to the shelter on the north side'') and are read as requests for shelter B: the explicit shelter rule refuses the names it knows (A, C, D, a gymnasium) and defaults to B for any other. In a world with one shelter the robot would deliver the right supply to the only shelter, but the operator asked for somewhere else, so the read-back confirmation of the deployment note is required, not optional. The fifth, ``record the number of water bottles delivered to shelter B'', is accepted as a request for water: a lexical trap that the closed-set question cannot tell from an order. Legitimate requests are read well (\IGAccepted{}; speech-recognition style \IGAsr{}, English \IGEnglish{}); the three falsely refused texts (``\texttt{B} first-aid set'' with no shelter word, a kana spelling of the first-aid box, and a negation that names the other supply) cost a question, not a delivery. Related work on voice-controlled embodied agents simulates speech-recognition errors on safety benchmarks and finds that they can make harmful instructions accepted~\cite{jia2026mishear}, and language models degrade on noisy instructions~\cite{wang2024noisyinstructions}; our texts are typed imitations of recogniser output and the failure we found is of a different kind, an unknown destination. Cross-lingual language-to-action grounding for robots has been shown for more than 140 languages~\cite{nwankwo2025reli} and vision-language-action models degrade outside English; Japanese has little standardised orthography, and speech-recognition evaluation therefore needs lenient scoring of valid spelling variants~\cite{karita2023lenient}, a block (kanji, hiragana, katakana and romaji variants of each destination name, with and without particles) that the next text set should contain; the quality of the test items itself matters, as the refinement of Japanese VLM benchmarks shows~\cite{sugiura2026jammeval}, and the 136 texts here were written by one assistant. Speech interfaces can lose task-critical details before execution~\cite{chiba2026talk2agent}, and flagging uncertain spans helps listeners catch recognition errors~\cite{nowrin2024asrconfidence}. The parser was not tuned on this set. Limits: 136 texts written by the same assistant that wrote the prompts, typed imitations of recogniser output, Japanese and English only. A change that asks for the destination as a closed-set question would be a new parser and needs a new fresh set, not a re-reading of this one.
